\section{Colour interpolation}
\label{sec:color}
\here{9}
Colour is interpolated in four places: across a triangle (barycentric), between shading models (Gouraud vs Phong),
between colour spaces (sRGB vs linear), and over time (sky and fog presets over the day). Each has a naive version
that fails visibly.
\begin{equation}
c_{\text{lin}}=\begin{cases}c_s/12.92 & c_s\le0.04045\\ \big(\frac{c_s+0.055}{1.055}\big)^{2.4} & \text{else}\end{cases},\qquad
\operatorname{mix}_{\log}(a,b,t)=\exp\big((1-t)\ln a+t\ln b\big)
\label{eq:srgb}
\end{equation}

\begin{figure}[htbp]
\centering
{\footnotesize\textit{A. The worked triangle}}\\[0.5mm]
\wbblock{Color_Interp}{verts}\\[2mm]
\begin{minipage}[t]{0.6\linewidth}\centering{\footnotesize\textit{B. Barycentric weights}}\\[0.5mm]
\wbblock{Color_Interp}{bary}\end{minipage}\hfill
\begin{minipage}[t]{0.38\linewidth}\centering{\footnotesize\textit{D. Texture $u$: affine vs perspective-correct}}\\[0.5mm]
\wbblock{Color_Interp}{uv}\end{minipage}\\[2mm]
{\footnotesize\textit{C. Gouraud vs Phong at the same points}}\\[0.5mm]
\wbblock{Color_Interp}{gp}\\[2mm]
{\footnotesize\textit{E. sRGB vs linear blending of colours A and B}}\\[0.5mm]
\wbblock{Color_Interp}{srgb}\\[2mm]
{\footnotesize\textit{F. Sky zenith colour over the day}}\\[0.5mm]
\wbblock{Color_Interp}{sky}\\[1mm]
{\footnotesize\wbkey}
\caption{Workbook sheet \code{Color\_Interp}: barycentric weights at sample points of one triangle, Gouraud vs Phong
values at the same points, affine vs perspective-correct texture coordinates with $w$, sRGB vs linear blends of two
colours, and the sky colours tabulated over the day. \fullsheet{Color_Interp}}
\label{fig:wb-color}
\end{figure}

\begin{table}[htbp]
\caption{Colour interpolation: the naive method, how it fails, and the renderer's method.}
\centering\small
\begin{tabularx}{\linewidth}{@{}lXXl@{}}
\toprule
Interpolation & Naive version and its failure & What the renderer does & Evidence\\\midrule
Across a triangle & affine in screen space: textures bend & perspective-correct ($a/w$, $1/w$) & Fig.~\ref{fig:raster}b\\
Lighting & Gouraud: highlight between vertices lost, Mach bands & Phong per fragment & Figs.~\ref{fig:shading}, \ref{fig:gp}\\
Colour space & blending sRGB values: mid-tones too dark, e.g.\ 50\,\% red/green gives a muddy (0.5, 0.5, 0) & material colours converted to linear, lit and blended in linear HDR, encoded once at the end & Fig.~\ref{fig:wb-color}\\
Sky over the day & linear mix of night and day presets: the night tint is swamped by the day value until very late & presets blended in log space (Eq.~\ref{eq:srgb}): equal steps are equal ratios, which the eye sees as even & Figs.~\ref{fig:hero}, \ref{fig:daycycle}\\
Heat to false colour & one linear ramp black$\to$white hides small heat differences & 7 colour stops, linear between neighbours (ironbow, rainbow), or a power curve (white-hot) & Fig.~\ref{fig:palettes}\\
\bottomrule
\end{tabularx}
\end{table}

\begin{figure}[htbp]
\centering
\plot{p_palettes}{\linewidth}\\[1.5mm]
\begin{subfigure}{0.325\linewidth}\img{figures/img/annot/pal_ironbow.jpg}{\linewidth}\caption{Ironbow.}\end{subfigure}\hfill
\begin{subfigure}{0.325\linewidth}\img{figures/img/annot/pal_whitehot.jpg}{\linewidth}\caption{White-hot.}\end{subfigure}\hfill
\begin{subfigure}{0.325\linewidth}\img{figures/img/annot/pal_rainbow.jpg}{\linewidth}\caption{Rainbow.}\end{subfigure}
\caption{Colour interpolation in the thermal view: the heat of each pixel ($t\in[0,1]$) is looked up in a palette that
interpolates linearly between colour stops (top, from \code{thermalPalette} in the shader, shown sRGB-encoded); (a--c)
the same night frame through each palette (\key{T} cycles them).}
\label{fig:palettes}
\end{figure}

\textbf{Limitation.} Log-space blending is undefined for a zero channel: every preset must stay strictly positive,
and a preset with a near-zero channel makes the blend linger near black for most of the transition. The final image
is stored in 8 bits per channel, and the only defence against banding in dark gradients is the film grain added in
the last pass.
